\section{Aproximação e acoplamento ou atracamento}
\label{sec_rev_bib_rvdb}

Conforme \cite{nasa1stDocking}, o primeiro encontro (\textit{rendezvous}) com acoplamento (\textit{docking}) entre duas espaçonaves ocorreu em $1966$, quando a tripulação do programa \emph{Gemini} realizou manualmente a aproximação e o acoplamento ao veículo não tripulado \emph{Agena}.

Como explicado por \cite{nasaAutmaticDocking1967}, pouco depois, em $1967$, foi demonstrada a viabilidade do acoplamento automático em órbita, com as missões soviéticas \emph{Kosmos} 186 e 188.

Os estudos de \cite{goodman2011rendezvous} abragem a experiência acumulada em missões \textit{Gemini} e \textit{Apollo} ($1968$--$1972$) que consolidaram os procedimentos de encontro e acoplamento, incluindo as fases de aproximação progressiva, pontos de espera e decisões operacionais baseadas na avaliação contínua do estado relativo.

O trabalho de \cite{goodman2011rendezvous} apresentam os conceitos que posteriormente foram aplicados em missões como \emph{Skylab} ($1973$--$1974$) e \emph{Apollo--Soyuz}, evidenciando a maturidade das operações de \gls{rvdb} como capacidade habilitadora para arquiteturas mais complexas.







% \subsection{Desenvolvimento das operações de acoplamento}
\label{subsec_desenvolvimento_acoplamento}

Após as primeiras demonstrações de \textit{rendezvous} orbital, o desenvolvimento das operações de acoplamento passou a concentrar esforços na criação de sistemas capazes de realizar a aproximação final e o contato mecânico entre veículos espaciais de forma segura e repetível. 
Esse avanço envolveu tanto a evolução de mecanismos de acoplamento quanto o desenvolvimento de sistemas de navegação relativa, guiamento e controle capazes de operar em distâncias progressivamente menores entre as espaçonaves.

Nas primeiras aplicações, as operações de acoplamento dependiam fortemente da atuação de pilotos humanos, responsáveis por interpretar informações de navegação relativa e conduzir manualmente as fases finais da aproximação. 
Com o avanço da tecnologia de sensores e de sistemas embarcados, passaram a ser desenvolvidas arquiteturas capazes de automatizar parte dessas operações, reduzindo a carga operacional sobre os astronautas e aumentando a previsibilidade do processo de aproximação.

Um dos desafios associados a essa evolução consiste em garantir a estimativa precisa da posição e orientação relativas entre os veículos durante a fase final de aproximação. 
Nesse contexto, diferentes soluções foram propostas ao longo do tempo, incluindo sistemas baseados em sensores ópticos, sensores a laser e técnicas de fusão de dados para estimar o estado relativo entre perseguidor e alvo.

Estudos iniciais sobre acoplamento automático entre espaçonaves investigaram arquiteturas de sensoriamento e controle capazes de conduzir o processo sem intervenção humana direta. 
Um exemplo é o trabalho apresentado por \cite{micheal1981}, no qual foi analisado um cenário de acoplamento associado ao projeto \emph{Teleoperator Retrieval System (TRS)}. 
Nesse estudo, foram avaliados métodos para estimativa de posição e atitude relativas utilizando sensores ópticos e sistemas de varredura a laser, evidenciando a importância da integração entre sensoriamento e algoritmos de controle para viabilizar operações autônomas de acoplamento.

Com a evolução das missões de logística orbital e de suporte a estações espaciais, sistemas de acoplamento altamente automatizados passaram a ser empregados em veículos de serviço. 
Um exemplo representativo é o \emph{Automated Transfer Vehicle} (ATV), desenvolvido pela \gls{esa} para missões de reabastecimento da \gls{iss}. 
O ATV realizava operações de \gls{rvdb} utilizando sistemas automáticos de navegação relativa, integrando sensores, câmeras e sistemas de guiamento para ajustar continuamente sua trajetória durante a fase final de aproximação \cite{esaATVRendezvous}.

Durante a aproximação final, o veículo utilizava sensores a laser para estimar a distância e o alinhamento em relação à estação espacial, permitindo realizar correções de trajetória em tempo real. 
A automação dessas operações reduziu significativamente a necessidade de intervenção humana direta e demonstrou a viabilidade de sistemas de \gls{rvdb} altamente automatizados em missões operacionais.

% \subsection{Operações de proximidade em estações espaciais}
\label{subsec_operacoes_proximidade_estacoes}

A evolução das operações de \gls{rvdb} possibilitou a realização de missões logísticas e de manutenção em estações espaciais, nas quais veículos visitantes realizam aproximações controladas para transporte de tripulação, suprimentos e equipamentos científicos. 
Essas operações exigem elevada precisão na navegação relativa e no controle de atitude durante as fases finais de aproximação, uma vez que pequenas incertezas na estimativa de posição ou orientação podem comprometer a segurança do encontro orbital.

A Estação Espacial Internacional (\gls{iss}) representa um dos principais exemplos de infraestrutura orbital que depende de operações recorrentes de proximidade. 
Desde o início de sua montagem, múltiplos veículos visitantes passaram a realizar missões de reabastecimento e transporte de tripulação, exigindo o desenvolvimento de diferentes estratégias de aproximação e acoplamento.

Além das operações de acoplamento direto (\emph{docking}), algumas missões utilizam o processo de atracação assistida (\emph{berthing}), no qual o veículo visitante é capturado por um \gls{mr} e posteriormente conduzido para acoplar-se à estação espacial. 
Nesse tipo de operação, o veículo aproxima-se da estação até uma posição de captura segura, permanecendo estacionário enquanto o manipulador realiza a captura e conduz o veículo até a interface de conexão.

Entre os exemplos de veículos que utilizam esse tipo de operação está a espaçonave cargueira \emph{Cygnus}, desenvolvida pela empresa \emph{Northrop Grumman} para transportar suprimentos e experimentos científicos para a \gls{iss}. 
Durante suas missões, a \emph{Cygnus} realiza a aproximação até uma posição de captura, sendo então agarrada pelo \gls{mr} \emph{Canadarm2}, operado pelos astronautas a bordo da estação, que finalizam o processo de atracação conectando o veículo a um dos módulos da estação \cite{nasaCygnus}.

De forma semelhante, o veículo japonês \emph{H-II Transfer Vehicle} (HTV), desenvolvido pela \gls{jaxa}, também foi projetado para realizar operações de \emph{berthing}. 
Após realizar a aproximação controlada à estação espacial, o HTV era capturado pelo \gls{mr} e posteriormente conectado à estrutura da \gls{iss}, permitindo o acesso às cargas transportadas \cite{nasaHTV}.

O \gls{mr} \emph{Canadarm2}, desenvolvido pela \gls{csa}, desempenha papel central nessas operações. 
Com múltiplos graus de liberdade e capacidade de ser operado tanto por astronautas quanto por controladores em solo, esse sistema permite capturar veículos visitantes, reposicionar módulos da estação e realizar tarefas de manutenção externa \cite{canadarm2}.


% \subsection{Missões recentes e novos desafios}
\label{subsec_missoes_recentes_desafios}

Nos últimos anos, as operações de \gls{rvdb} passaram a ocorrer em um contexto mais amplo, caracterizado pela crescente participação de empresas privadas, pelo desenvolvimento de serviços em órbita e pela preparação de missões de exploração além da órbita baixa terrestre. 
Essas mudanças ampliaram o número de aplicações das operações de aproximação e acoplamento, ao mesmo tempo em que introduziram novos requisitos tecnológicos e operacionais.

Um exemplo desse cenário é a expansão de veículos comerciais destinados ao transporte de carga e tripulação para a \gls{iss}. 
A espaçonave \emph{Dragon}, desenvolvida pela \emph{SpaceX}, tornou-se um dos principais veículos logísticos da estação espacial, sendo responsável pelo transporte de suprimentos e pelo retorno de experimentos científicos para a Terra \cite{bell2018,nasaDragonReturn}. 
Em suas missões, a \emph{Dragon} realiza operações de aproximação utilizando sistemas de navegação relativa baseados em sensores como LiDAR e GPS, que permitem estimar a posição e orientação em relação à estação durante as fases finais de aproximação \cite{nasaspaceflightDragonISS}.

A evolução desses sistemas levou ao desenvolvimento de veículos com maior grau de autonomia. 
A versão mais recente da espaçonave, \emph{Dragon 2}, incorpora capacidades de navegação e acoplamento autônomos, permitindo a realização de \emph{docking} direto com a estação espacial sem a necessidade do manipulador robótico \emph{Canadarm2} \cite{nasaStarshipDocking}. 
Esse tipo de arquitetura reduz a dependência de intervenção humana durante as fases críticas da missão e aumenta a robustez operacional das operações de \gls{rvdb}.

Outro campo de expansão recente das operações de proximidade envolve missões de manutenção e extensão da vida útil de satélites em órbita. 
O \gls{mev}, desenvolvido pela \emph{Northrop Grumman}, demonstrou a viabilidade de operações de acoplamento com satélites previamente lançados, permitindo assumir funções de controle de atitude e manutenção orbital do veículo atendido \cite{mev1}. 
A primeira missão operacional ocorreu em 2020, quando o \emph{MEV-1} acoplou-se ao satélite \emph{Intelsat 901}, prolongando sua vida útil operacional \cite{nytMEV1}.

Além das aplicações em órbita terrestre, futuras arquiteturas de exploração espacial também dependem de capacidades avançadas de \gls{rvdb}. 
O programa \emph{Artemis}, liderado pela \gls{nasa}, prevê a construção de infraestrutura orbital no ambiente cislunar, incluindo a estação \emph{Gateway} \cite{nasaArtemis}. 
Nesse contexto, operações de aproximação e acoplamento devem ser realizadas com maior grau de autonomia, uma vez que atrasos de comunicação entre a Terra e o ambiente lunar podem limitar a supervisão contínua por operadores em solo \cite{keeports2006}.

Esses cenários reforçam a importância do desenvolvimento de sistemas de navegação relativa, sensoriamento e controle capazes de operar de forma robusta em diferentes regimes orbitais. 
O uso de sensores como LiDAR, câmeras estereoscópicas e algoritmos avançados de processamento de dados permite estimar estados relativos e apoiar o planejamento de trajetórias de aproximação com menor dependência de intervenção humana direta \cite{farkasvolgyi2023,ibmComputerVision}.




